\chapter{机械臂动力学}
\label{ch:arm-dyn}

% ════════════════════════════════════════════════════════════════
%  机械臂建模、规划与控制部 · 第 2 章。定基串联臂动力学。
%  设计纪律（融入设计_机械臂新部.md §0 复用纪律）：
%   - 理论无教材抽取源，按编写规范 §一.9 用网络重建到「复现级」，§一.8 自包含写入正文；
%     存疑公式打 \rebuilt。出处 \cite featherstone2008rbda / spong2006robot / lynch2017modern
%     / khatib1987operational。
%   - P4 \cref{sec:ctrl-robot-dyn} 已给方程 \cref{eq:ctrl-manip} 与结构性质
%     \cref{prop:ctrl-robot-props}（结论）；本章给「完整推导来历」，绝不重列结论、只 \cref。
%   - RNEA 复用 P7 \cref{sec:qk-plucker} 的 Plücker 空间向量代数（记号不重建）。
%   - §2.6 条件数 = Robot吸收_arm机械臂运控.md §三根因 2/3 的最深实证（cond 527-1625 /
%     正则破坏投影性 / 轻腕欠采样）。实践接 \cref{ch:arm-prac}。
%  教学层遵循 v5.0：动机→反面→理论→案例→陷阱→练习 + 符号表/速查/故障排查。
% ════════════════════════════════════════════════════════════════

\begin{practice}[前置自测]
开始之前，先试答下列问题。若已能流畅作答，可只速读本章的速查卡与符号表；若卡住，正是本章要为你解决的。
\begin{enumerate}
  \item 上一章（\cref{ch:arm-kin}）把“关节角 $\leftrightarrow$ 末端位姿/速度”讲透了，纯运动学已经能让末端\emph{走到}任何可达点。那么为什么还需要\emph{动力学}？请举出至少三件“只有运动学、没有动力学就做不成”的事。
  \item 同一个 $n$ 自由度机械臂，它的运动方程 $\mathbf{M}(\mathbf{q})\ddot{\mathbf{q}}+\mathbf{C}(\mathbf{q},\dot{\mathbf{q}})\dot{\mathbf{q}}+\mathbf{g}(\mathbf{q})=\boldsymbol{\tau}$ 既可以从“能量”（拉格朗日）推出，也可以从“逐杆受力”（牛顿--欧拉）推出。这两条路得到的是\emph{同一个}方程吗？各自的优势是什么——一个适合\emph{看清结构}，另一个适合\emph{快速计算}，分别是哪一个？
  \item 给定 $(\mathbf{q},\dot{\mathbf{q}},\ddot{\mathbf{q}})$ 反求所需力矩 $\boldsymbol{\tau}$（\emph{逆动力学}），朴素地按定义组装 $\mathbf{M},\mathbf{C},\mathbf{g}$ 再相乘需要多少次运算（按 $n$ 计）？有没有一种递推算法能把它降到 $O(n)$、且\emph{根本不显式形成} $\mathbf{M}$？
  \item 质量阵 $\mathbf{M}(\mathbf{q})$ 总是对称正定（\cref{prop:ctrl-robot-props}），所以总是“可逆”的。可是“可逆”和“求逆数值稳定”是两回事。一台“重肩肘 + 轻腕”的机械臂，腕部惯量只有 $2\times10^{-4}$，肩肘惯量却大三个数量级——这会让 $\mathbf{M}$ 的\emph{条件数}发生什么？它进而会让哪些依赖 $\mathbf{M}^{-1}$ 的控制律（如操作空间控制）\emph{失效}？
  \item 操作空间控制要把关节力矩 $\boldsymbol{\tau}$ 与末端力 $\mathbf{F}$ 对应起来。除了\emph{运动学}对偶 $\boldsymbol{\tau}=\mathbf{J}^\top\mathbf{F}$，为什么还要一个“任务空间惯量” $\boldsymbol{\Lambda}=(\mathbf{J}\mathbf{M}^{-1}\mathbf{J}^\top)^{-1}$？它解决了纯运动学映射解决不了的什么问题？
\end{enumerate}
\end{practice}

\section*{本章目标}
学完本章，你应当能够：
\begin{enumerate}
  \item \textbf{说清}机械臂为何需要动力学——力矩可行性、接触/力控、前馈跟踪三条动机，并理解“运动学规划的轨迹未必是力矩可执行的”（\cref{sec:ad-why}）；
  \item \textbf{完整推导}Euler--Lagrange 法如何从动能/势能得到 $\mathbf{M}(\mathbf{q})\ddot{\mathbf{q}}+\mathbf{C}\dot{\mathbf{q}}+\mathbf{g}=\boldsymbol{\tau}$，并用 Christoffel 符号写出科氏阵 $\mathbf{C}$、看清它与质量阵导数的关系（\cref{sec:ad-lagrange}）；
  \item \textbf{推导并实现}牛顿--欧拉递推（RNEA）：前向传播速度/加速度、后向传播力/力矩，得到 $O(n)$ 的逆动力学，并理解它与 Lagrange 法的等价（\cref{sec:ad-rnea}）；
  \item \textbf{用}复合刚体算法（CRBA）算质量阵 $\mathbf{M}(\mathbf{q})$，并理解 $\mathbf{M}\succ0$、$\dot{\mathbf{M}}-2\mathbf{C}$ 斜对称（无源性）的算法级来历（\cref{sec:ad-crba}）；
  \item \textbf{建立}操作空间动力学：任务空间惯量 $\boldsymbol{\Lambda}=(\mathbf{J}\mathbf{M}^{-1}\mathbf{J}^\top)^{-1}$、动力学一致广义逆 $\bar{\mathbf{J}}=\mathbf{M}^{-1}\mathbf{J}^\top\boldsymbol{\Lambda}$，为后续 \cref{ch:arm-ctrl} 的操作空间控制铺路（\cref{sec:ad-osd}）；
  \item \textbf{用第一性原理（条件数）判断算法可行性}：理解“重肩肘 + 轻腕”异质惯量臂的固有病态质量阵如何让操作空间控制收不住末端、如何让轻腕在控制频率下欠采样发散（\cref{sec:ad-conditioning}）；
  \item \textbf{查用}动力学计算工具（Pinocchio \texttt{rnea}/\allowbreak\texttt{crba}/\allowbreak\texttt{computeGeneralizedGravity}）的用途映射（\cref{sec:ad-practice}）。
\end{enumerate}

\subsection*{本章知识导航}
本章是“\emph{从关节力矩到关节加速度}”的桥：上一章（\cref{ch:arm-kin}）把几何讲透，本章补上\emph{物理}。主线有两条互补的推导路径——\emph{能量观}（Euler--Lagrange，看清结构）与\emph{递推观}（牛顿--欧拉 RNEA，$O(n)$ 快算），它们殊途同归到\emph{同一个}运动方程。理解了方程的来历与结构性质，才能在 \cref{ch:arm-traj} 把“力矩约束”塞进时间参数化、在 \cref{ch:arm-ctrl} 写出计算力矩/操作空间/阻抗控制。

\begin{center}
\begin{forest} navtreeLR
[{机械臂动力学\\$\mathbf{M}\ddot{\mathbf{q}}+\mathbf{C}\dot{\mathbf{q}}+\mathbf{g}=\boldsymbol{\tau}$}
  [{为何动力学\\力矩/接触/前馈} ]
  [{能量观\\Euler--Lagrange}
    [{Christoffel $\to\mathbf{C}$}] ]
  [{递推观\\牛顿--欧拉 RNEA}
    [{$O(n)$ 逆动力学}] ]
  [{质量阵\\CRBA + 结构性质}
    [{$\mathbf{M}\succ0$/斜对称}] ]
  [{操作空间动力学\\$\boldsymbol{\Lambda},\bar{\mathbf{J}}$}
    [{动力学一致逆}] ]
  [{条件数（实证）\\病态 $\mathbf{M}$}
    [{OSC 失效/轻腕发散}] ]
]
\end{forest}
\end{center}

\subsection*{前置桥接}
本章重度复用本书已建的地基，遇到时一律交叉引用回去、不重证：
\begin{itemize}
  \item \textbf{P0 刚体与李群}：位姿 $\mathbf{T}\in\SE(3)$（\cref{ch:rigid_body}、\cref{def:se3-rep}）、伴随变换 $\mathrm{Ad}$（\cref{eq:Ad}）、末端速度运动学 $\boldsymbol{\varpi}=\boldsymbol{\mathcal{J}}\dot{\boldsymbol{\xi}}$（\cref{thm:se3-kinematics}）、反对称恒等式 $(\mathbf{R}\mathbf{a})^\wedge=\mathbf{R}\mathbf{a}^\wedge\mathbf{R}^\top$（\cref{eq:rb-id5}）。
  \item \textbf{P7 空间向量代数}：RNEA 直接借用 \cref{sec:qk-plucker} 已建的 Plücker 六维空间向量（运动/力向量、空间变换 $\mathbf{X}$、空间叉乘、空间惯量），本章不重建这套记号，只把它用到\emph{定基串联臂}上。
  \item \textbf{P4 控制导论}：\cref{sec:ctrl-robot-dyn} 已经\emph{给出}方程 \cref{eq:ctrl-manip} 与两条结构性质（\cref{prop:ctrl-robot-props}）作为控制的起点——那里是“结论”，\textbf{本章是它的“完整推导来历”}。读者可对照：P4 直接拿来用，本章把每一项怎么来的讲清。
  \item \textbf{上一章雅可比}：操作空间动力学（\cref{sec:ad-osd}）用到几何雅可比 $\mathbf{J}$ 与静力学对偶 $\boldsymbol{\tau}=\mathbf{J}^\top\mathbf{F}$（\cref{sec:ak-jacobian} 已建）。
\end{itemize}

% ════════════════════════════════════════════════════════════════
\section{为何需要动力学：运动学不够用}
\label{sec:ad-why}

\paragraph{动机：能“走到”不等于能“走好”、更不等于能“接触”。}
上一章的纯运动学已经强大：给定关节角就知道末端在哪（正运动学），给定末端目标就反求关节角（逆运动学），给定关节速度就知道末端怎么动（雅可比）。在“自由空间、慢速、无接触”的世界里，运动学几乎够用——许多工业“示教--再现”就是纯运动学的位置回放。但只要踏出这个温室，三类问题立刻把动力学逼出来。

\begin{enumerate}
  \item \textbf{力矩可行性（能不能驱动）}。运动学规划出一条几何路径 $\mathbf{q}(s)$，再配上时间得到 $\mathbf{q}(t)$，于是有了 $\dot{\mathbf{q}},\ddot{\mathbf{q}}$。可\emph{电机能不能提供这条轨迹所需的力矩}？同一个关节加速度 $\ddot{\mathbf{q}}$，在不同位形下需要的力矩\emph{天差地别}——因为有效惯量 $\mathbf{M}(\mathbf{q})$ 随位形变。一个“限加速度”看似安全的轨迹，可能在某个位形要求 $>26\,\mathrm{N\!\cdot\! m}$ 而真机给不出；反过来，它也可能在大半段路上只用到驱动力的零头、白白浪费。\textbf{要回答“这条轨迹力矩可行吗”，必须有 $\boldsymbol{\tau}=\mathbf{M}\ddot{\mathbf{q}}+\mathbf{C}\dot{\mathbf{q}}+\mathbf{g}$}。这正是 \cref{ch:arm-traj} 力矩约束时间参数化的根。
  \item \textbf{接触与力控（出多大力）}。打磨、插轴、开门、与人协作——任务的核心是\emph{力}而非位置。要让末端“顶住墙面恒力 $15\,\mathrm{N}$”或“被推就柔顺退让”，控制器必须把“期望的力/柔顺”换算成“该发的关节力矩”。这个换算的桥就是动力学：静力学对偶 $\boldsymbol{\tau}=\mathbf{J}^\top\mathbf{F}$ 把末端力折成关节力矩，重力补偿 $\boldsymbol{\tau}=\mathbf{g}(\mathbf{q})$ 让臂“失重悬停”从而能呈现纯净的弹簧律。\textbf{没有动力学，力控无从谈起}（\cref{ch:arm-ctrl}）。
  \item \textbf{高速/高精度跟踪（前馈抵消非线性）}。要让末端\emph{精确}跟踪一条时变轨迹，纯反馈（PD）会因为机器人的强非线性、强耦合而滞后。\textbf{计算力矩控制}用模型\emph{前馈}抵消非线性，把闭环化为线性解耦误差动态——而那个前馈项 $\mathbf{M}\ddot{\mathbf{q}}_d+\mathbf{C}\dot{\mathbf{q}}+\mathbf{g}$ 就是逆动力学。
\end{enumerate}

\begin{insight}[运动学规划的轨迹未必是力矩可执行的]
\label{ins:ad-kin-not-enough}
这是贯穿本部的一条线：\textbf{运动学（\cref{ch:arm-kin}）回答“走哪条路、各点速度多少”，动力学回答“这条路要花多大力、能不能花得起”}。把“限加速度盒”当作力矩约束的代理是一种\emph{常见且危险}的近似——因为有效惯量 $\mathbf{M}(\mathbf{q})$ 与位形强相关，固定加速度盒在不同位形对应的力矩需求可以差好几倍（实证见 \cref{ch:arm-prac} 的力矩约束 TOPP-RA：把加速度约束换成真力矩约束后，轨迹快了 $4.6\times$ 还把驱动力用满）。一句话：\emph{动力学是“运动学画的饼”能否落地的可行性检验器}。
\end{insight}

\paragraph{本章要建立什么。}
我们要从两条互补的路推出同一个运动方程，并把它的\emph{结构}与\emph{算法}讲透：
\begin{equation}
  \boxed{\;\mathbf{M}(\mathbf{q})\ddot{\mathbf{q}}+\mathbf{C}(\mathbf{q},\dot{\mathbf{q}})\dot{\mathbf{q}}+\mathbf{g}(\mathbf{q})=\boldsymbol{\tau}+\textstyle\sum_k\mathbf{J}_k^\top\mathbf{F}_k\;}
  \label{eq:ad-manip}
\end{equation}
这与 P4 的 \cref{eq:ctrl-manip} 是\emph{同一个方程}（记号完全一致）。区别在视角：P4 把它当控制的\emph{起点}直接用；本章把每一项的来历、每个结构性质的根、以及怎么\emph{高效计算}它讲清。各符号：$\mathbf{q}\in\mathbb{R}^n$ 关节角，$\mathbf{M}(\mathbf{q})\succ0$ 质量（惯性）阵，$\mathbf{C}(\mathbf{q},\dot{\mathbf{q}})\dot{\mathbf{q}}$ 科氏/离心项，$\mathbf{g}(\mathbf{q})$ 重力力矩，$\boldsymbol{\tau}$ 关节驱动力矩，$\mathbf{F}_k$ 第 $k$ 个接触/外力（经其雅可比 $\mathbf{J}_k$ 折成广义力）。

\paragraph{两个对偶问题：逆动力学 vs 正动力学。}
动力学有两个互为表里的计算问题，全章反复用到，先厘清：
\begin{itemize}
  \item \textbf{逆动力学}（inverse dynamics, ID）：已知 $(\mathbf{q},\dot{\mathbf{q}},\ddot{\mathbf{q}})$，求所需力矩 $\boldsymbol{\tau}=\mathrm{ID}(\mathbf{q},\dot{\mathbf{q}},\ddot{\mathbf{q}})$。这是\emph{控制/规划}最常用的方向（前馈、力矩约束、计算力矩都靠它）。RNEA（\cref{sec:ad-rnea}）以 $O(n)$ 解它。
  \item \textbf{正动力学}（forward dynamics, FD）：已知 $(\mathbf{q},\dot{\mathbf{q}},\boldsymbol{\tau})$，求加速度 $\ddot{\mathbf{q}}=\mathbf{M}^{-1}(\boldsymbol{\tau}-\mathbf{C}\dot{\mathbf{q}}-\mathbf{g})$。这是\emph{仿真}要算的方向（物理引擎积分它）。它需要 $\mathbf{M}^{-1}$，于是要么用 CRBA（\cref{sec:ad-crba}）显式算 $\mathbf{M}$ 再解线性系统，要么用关节体算法 ABA（本章只点名，详见 \cite{featherstone2008rbda}）以 $O(n)$ 直接得 $\ddot{\mathbf{q}}$。
\end{itemize}
本章主攻 ID 与 $\mathbf{M}$ 的计算（控制最需要的），FD 的 $O(n)$ 算法（ABA）作为延伸留给文献。

% ════════════════════════════════════════════════════════════════
\section{Euler--Lagrange 法：从能量看清结构}
\label{sec:ad-lagrange}

\paragraph{动机：为什么先讲能量法。}
牛顿第二定律“逐杆受力分析”当然能列方程，但对多连杆耦合系统，逐杆引入内部约束力（关节反力）会让推导陷入泥潭——而这些内部力\emph{最终都消去了}，是无用功。\textbf{拉格朗日法的高明在于：只用标量能量（动能、势能）和广义坐标，内部约束力自动消去，且推出的方程\emph{结构清晰}}——质量阵、科氏阵、重力项各就各位。本节用它\emph{看清结构}（结构性质的证明最自然）；下一节再用牛顿--欧拉\emph{快速计算}。两者结果相同（\cref{ins:ad-equiv}）。

\paragraph{拉格朗日量与 Euler--Lagrange 方程。}
取广义坐标为关节角 $\mathbf{q}=[q_1,\dots,q_n]^\top$。系统\emph{拉格朗日量} $L=K-P$（动能减势能）。\emph{Euler--Lagrange 方程}给出广义力（这里即关节力矩 $\boldsymbol{\tau}$）与运动的关系\cite{spong2006robot,lynch2017modern}：
\begin{equation}
  \boxed{\;\frac{\mathrm{d}}{\mathrm{d}t}\frac{\partial L}{\partial\dot{q}_i}-\frac{\partial L}{\partial q_i}=\tau_i,\qquad i=1,\dots,n.\;}
  \label{eq:ad-el}
\end{equation}
\cref{eq:ad-el} 是\emph{解析力学}的基本定理（其证明属变分原理范畴，是本书的假定背景，见 \cite{tedrake_underactuated}），我们直接\emph{用}它来推机械臂方程。核心工作是把 $K,P$ 用 $\mathbf{q},\dot{\mathbf{q}}$ 表出。

\paragraph{动能与质量阵的来历。}
关键观察：\textbf{机械臂动能是关节速度的二次型，其系数矩阵就是质量阵 $\mathbf{M}(\mathbf{q})$}。逐杆累加：第 $i$ 杆质心线速度 $\mathbf{v}_{c_i}$ 与角速度 $\boldsymbol{\omega}_i$ 都\emph{线性}依赖 $\dot{\mathbf{q}}$——由上一章雅可比，存在质心雅可比 $\mathbf{J}_{v_i}(\mathbf{q}),\ \mathbf{J}_{\omega_i}(\mathbf{q})$ 使
\begin{equation}
  \mathbf{v}_{c_i}=\mathbf{J}_{v_i}(\mathbf{q})\,\dot{\mathbf{q}},\qquad \boldsymbol{\omega}_i=\mathbf{J}_{\omega_i}(\mathbf{q})\,\dot{\mathbf{q}}.
  \label{eq:ad-link-vel}
\end{equation}
第 $i$ 杆动能 = 平动 $\tfrac12 m_i\|\mathbf{v}_{c_i}\|^2$ 加转动 $\tfrac12\boldsymbol{\omega}_i^\top\mathbf{I}_{c_i}\boldsymbol{\omega}_i$（$m_i$ 质量、$\mathbf{I}_{c_i}$ 质心惯量张量，须在同一坐标系，记 $\mathbf{I}_{c_i}=\mathbf{R}_i\,{}^{c_i}\!\mathbf{I}_{c_i}\,\mathbf{R}_i^\top$ 由机体惯量旋到世界系）。全臂动能：
\begin{equation}
  K=\frac12\dot{\mathbf{q}}^\top\!\underbrace{\left[\sum_{i=1}^n\Big(m_i\mathbf{J}_{v_i}^\top\mathbf{J}_{v_i}+\mathbf{J}_{\omega_i}^\top\mathbf{I}_{c_i}\mathbf{J}_{\omega_i}\Big)\right]}_{\displaystyle\mathbf{M}(\mathbf{q})}\dot{\mathbf{q}}=\frac12\dot{\mathbf{q}}^\top\mathbf{M}(\mathbf{q})\dot{\mathbf{q}}.
  \label{eq:ad-kinetic}
\end{equation}

\begin{insight}[质量阵的两个性质从这里就读得出]
\label{ins:ad-M-spd}
从 \cref{eq:ad-kinetic} 的构造立刻看出 $\mathbf{M}(\mathbf{q})$ 的两条性质（对应 P4 \cref{prop:ctrl-robot-props} 性质 1）：(i) \textbf{对称}——每个被加项 $m_i\mathbf{J}_{v_i}^\top\mathbf{J}_{v_i}$ 与 $\mathbf{J}_{\omega_i}^\top\mathbf{I}_{c_i}\mathbf{J}_{\omega_i}$ 都是“$\mathbf{A}^\top(\cdot)\mathbf{A}$”形式、天然对称（$\mathbf{I}_{c_i}\succ0$ 对称）；(ii) \textbf{正定}——动能 $K=\tfrac12\dot{\mathbf{q}}^\top\mathbf{M}\dot{\mathbf{q}}>0$ 对任何 $\dot{\mathbf{q}}\neq\mathbf{0}$（只要每杆有质量、非退化位形），故 $\mathbf{M}\succ0$。正定意味着\emph{处处可逆}——但“可逆”不等于“求逆数值稳定”，这道伏笔将在 \cref{sec:ad-conditioning} 引爆。
\end{insight}

\paragraph{势能与重力项。}
势能只来自重力（弹性、关节弹簧暂不计）：$P(\mathbf{q})=\sum_i m_i\,\mathbf{g}_0^\top\mathbf{p}_{c_i}(\mathbf{q})$，其中 $\mathbf{g}_0$ 取\emph{重力加速度的相反矢量}（竖直向上、量值 $9.81$），如此 $P$ 随质心升高而增大、与“重力势能随高度增加”一致\footnote{此即 Modern Robotics \cite{lynch2017modern} 第 8 章的约定：把 $\mathbf{g}_0$ 写成向上矢量，势能项 $P=\sum_i m_i\mathbf{g}_0^\top\mathbf{p}_{c_i}$ 不带负号；等价写法是取向下的重力加速度 $\mathbf{a}_g$ 而令 $P=-\sum_i m_i\mathbf{a}_g^\top\mathbf{p}_{c_i}$。两者给出相同的 $\mathbf{g}(\mathbf{q})$。注意这与 \cref{alg:ad-rnea} 的重力技巧 $\mathbf{a}_0=-\mathbf{a}_g$（基座向上加速）正好用同一向上矢量。}；$\mathbf{p}_{c_i}$ 第 $i$ 杆质心位置（关节角的函数）。重力力矩定义为势能梯度：
\begin{equation}
  \mathbf{g}(\mathbf{q})\;\triangleq\;\frac{\partial P}{\partial\mathbf{q}}\;=\;\sum_{i=1}^n m_i\,\mathbf{J}_{v_i}^\top(\mathbf{q})\,\mathbf{g}_0,
  \label{eq:ad-gravity}
\end{equation}
最后一步用 $\partial\mathbf{p}_{c_i}/\partial\mathbf{q}=\mathbf{J}_{v_i}$（质心位置对关节角的雅可比即质心线速度雅可比）。这正是重力补偿控制 $\boldsymbol{\tau}=\mathbf{g}(\mathbf{q})$ 用 Pinocchio \texttt{computeGeneralizedGravity} 算的量（\cref{ch:arm-prac}）。

\paragraph{代入 Euler--Lagrange，凑出运动方程。}
$P$ 不含 $\dot{\mathbf{q}}$，故 $\partial L/\partial\dot{q}_i=\partial K/\partial\dot{q}_i=\sum_j M_{ij}\dot{q}_j$（用 $\mathbf{M}$ 对称）。逐项算 \cref{eq:ad-el}：

\begin{derivation}[从 Euler--Lagrange 凑出 $\mathbf{M}\ddot{\mathbf{q}}+\mathbf{C}\dot{\mathbf{q}}+\mathbf{g}=\boldsymbol{\tau}$]
第一项（对时间求导）：
\[
  \frac{\mathrm{d}}{\mathrm{d}t}\frac{\partial K}{\partial\dot{q}_i}=\frac{\mathrm{d}}{\mathrm{d}t}\sum_j M_{ij}\dot{q}_j=\sum_j M_{ij}\ddot{q}_j+\sum_j\dot{M}_{ij}\dot{q}_j=\sum_j M_{ij}\ddot{q}_j+\sum_{j,k}\frac{\partial M_{ij}}{\partial q_k}\dot{q}_k\dot{q}_j.
\]
第二项（对坐标求导，仅动能含 $q$ 在 $\mathbf{M}(\mathbf{q})$ 内）：$\dfrac{\partial K}{\partial q_i}=\dfrac12\sum_{j,k}\dfrac{\partial M_{jk}}{\partial q_i}\dot{q}_j\dot{q}_k$。势能项 $-\partial(-P)/\partial q_i=\partial P/\partial q_i=g_i$。三者并入 \cref{eq:ad-el}：
\[
  \sum_j M_{ij}\ddot{q}_j+\underbrace{\sum_{j,k}\Big(\frac{\partial M_{ij}}{\partial q_k}-\frac12\frac{\partial M_{jk}}{\partial q_i}\Big)\dot{q}_j\dot{q}_k}_{\text{科氏/离心}}+g_i=\tau_i.
\]
把中间二次速度项的系数对称化（利用对 $j,k$ 求和的对称性，把 $\partial M_{ij}/\partial q_k$ 与 $\partial M_{ik}/\partial q_j$ 平均），即得\emph{第一类 Christoffel 符号}
\begin{equation}
  c_{ijk}\;=\;\frac12\Big(\frac{\partial M_{ij}}{\partial q_k}+\frac{\partial M_{ik}}{\partial q_j}-\frac{\partial M_{jk}}{\partial q_i}\Big),
  \label{eq:ad-christoffel}
\end{equation}
于是科氏阵元素 $C_{ij}=\sum_k c_{ijk}\dot{q}_k$，整理成矩阵形式即得 \cref{eq:ad-manip}（无外力）。
\end{derivation}

\paragraph{科氏阵的 Christoffel 写法与“反对称”的由来。}
\cref{eq:ad-christoffel} 给出科氏阵的\emph{标准}选法
\begin{equation}
  C_{ij}(\mathbf{q},\dot{\mathbf{q}})=\sum_{k=1}^n c_{ijk}\dot{q}_k=\sum_{k=1}^n\frac12\Big(\frac{\partial M_{ij}}{\partial q_k}+\frac{\partial M_{ik}}{\partial q_j}-\frac{\partial M_{jk}}{\partial q_i}\Big)\dot{q}_k.
  \label{eq:ad-coriolis-matrix}
\end{equation}
注意 $\mathbf{C}$ \emph{不唯一}（满足 $\mathbf{C}\dot{\mathbf{q}}=$ 同一向量的 $\mathbf{C}$ 有无穷多个），但 Christoffel 选法有一条黄金性质：

\begin{proposition}[$\dot{\mathbf{M}}-2\mathbf{C}$ 斜对称（无源性的根）]
\label{prop:ad-skew}
取 \cref{eq:ad-coriolis-matrix} 的 Christoffel 科氏阵，则 $\mathbf{N}(\mathbf{q},\dot{\mathbf{q}})\triangleq\dot{\mathbf{M}}(\mathbf{q})-2\mathbf{C}(\mathbf{q},\dot{\mathbf{q}})$ 斜对称：$\mathbf{z}^\top\mathbf{N}\mathbf{z}=0,\ \forall\mathbf{z}$。
\end{proposition}
\begin{proof}\rebuilt{（据 Christoffel 定义直接验，参 \cite{spong2006robot}）}
逐元素：$\dot{M}_{ij}=\sum_k(\partial M_{ij}/\partial q_k)\dot{q}_k$，而 $2C_{ij}=\sum_k(\partial M_{ij}/\partial q_k+\partial M_{ik}/\partial q_j-\partial M_{jk}/\partial q_i)\dot{q}_k$。相减：
\[
  N_{ij}=\dot{M}_{ij}-2C_{ij}=\sum_k\Big(\frac{\partial M_{jk}}{\partial q_i}-\frac{\partial M_{ik}}{\partial q_j}\Big)\dot{q}_k.
\]
交换 $i\leftrightarrow j$：$N_{ji}=\sum_k(\partial M_{ik}/\partial q_j-\partial M_{jk}/\partial q_i)\dot{q}_k=-N_{ij}$（用 $\mathbf{M}$ 对称 $M_{jk}=M_{kj}$ 不改偏导记号）。故 $\mathbf{N}^\top=-\mathbf{N}$。
\end{proof}

\begin{insight}[斜对称即“内力不做功”——能量观的馈赠]
\label{ins:ad-passivity}
\cref{prop:ad-skew} 不是巧合，而是\emph{物理}的代数显影：惯性力与科氏/离心力\emph{既不产生也不耗散能量}。由动能 $K=\tfrac12\dot{\mathbf{q}}^\top\mathbf{M}\dot{\mathbf{q}}$，$\dot K=\dot{\mathbf{q}}^\top\mathbf{M}\ddot{\mathbf{q}}+\tfrac12\dot{\mathbf{q}}^\top\dot{\mathbf{M}}\dot{\mathbf{q}}$；代入 $\mathbf{M}\ddot{\mathbf{q}}=\boldsymbol{\tau}-\mathbf{C}\dot{\mathbf{q}}-\mathbf{g}$ 得 $\dot K=\dot{\mathbf{q}}^\top(\boldsymbol{\tau}-\mathbf{g})+\tfrac12\dot{\mathbf{q}}^\top(\dot{\mathbf{M}}-2\mathbf{C})\dot{\mathbf{q}}$，末项为零即“功率 = 外力做功率 − 重力做功率”，无凭空能量。\textbf{这条性质是机器人控制稳定性证明的“瑞士军刀”}：在 Lyapunov 求导中科氏项总与 $\tfrac12\dot{\mathbf{q}}^\top\dot{\mathbf{M}}\dot{\mathbf{q}}$ 配对消掉，于是\emph{不必知道 $\mathbf{C}$ 的具体形式}就能证稳定——PD+重力补偿（\cref{thm:ctrl-pd-grav}）正是靠它。这与 P4 \cref{prop:ctrl-robot-props} 的性质 2 是\emph{同一条}，本章给的是它\emph{从 Christoffel 定义出发}的算法级证明。
\end{insight}

\paragraph{第三条结构性质：惯性参数线性可参数化。}
还有一条对\emph{自适应控制}至关重要的性质（这里只点出、不展开，详见 \cite{spong2006robot}）：动力学对惯性参数 $\boldsymbol{\Theta}$（各杆质量、质心、惯量元素）\emph{线性}——存在回归阵 $\mathbf{Y}$ 使
\begin{equation}
  \mathbf{M}(\mathbf{q})\ddot{\mathbf{q}}+\mathbf{C}(\mathbf{q},\dot{\mathbf{q}})\dot{\mathbf{q}}+\mathbf{g}(\mathbf{q})=\mathbf{Y}(\mathbf{q},\dot{\mathbf{q}},\ddot{\mathbf{q}})\,\boldsymbol{\Theta}.
  \label{eq:ad-linear-param}
\end{equation}
这让“未知惯量在线辨识”成为可能——自适应控制器估计 $\boldsymbol{\Theta}$ 而非整条非线性方程。\cref{eq:ad-linear-param} 对应 P4 \cref{prop:ctrl-robot-props} 的性质 3。

\begin{pitfall}[概念误区：以为 $\mathbf{C}$ 是唯一的、或“$\mathbf{C}\dot{\mathbf{q}}$ 必须按 Christoffel 算”]
\label{pit:ad-coriolis-unique}
两个常见误解。其一，\textbf{科氏阵 $\mathbf{C}$ 不唯一}——满足 $\mathbf{C}(\mathbf{q},\dot{\mathbf{q}})\dot{\mathbf{q}}$ 等于那条二次速度向量的 $\mathbf{C}$ 有无穷多个（任何使 $(\mathbf{C}-\mathbf{C}')\dot{\mathbf{q}}=\mathbf{0}$ 的 $\mathbf{C}'$ 都行）。\cref{prop:ad-skew} 的斜对称\emph{只对 Christoffel 选法（\cref{eq:ad-coriolis-matrix}）成立}；若你从别处拿来一个 $\mathbf{C}$，先验证它的斜对称再用。其二，\textbf{实际计算逆动力学时几乎\emph{不}显式组装 $\mathbf{C}$}——RNEA（\cref{sec:ad-rnea}）一次递推就把 $\mathbf{C}\dot{\mathbf{q}}+\mathbf{g}$（甚至 $\mathbf{M}\ddot{\mathbf{q}}$）一起算出，根本不形成 Christoffel 符号那 $n^3$ 个偏导。Christoffel 写法的价值在\emph{看清结构与证性质}，不在\emph{算}。
\end{pitfall}

% ════════════════════════════════════════════════════════════════
\section{牛顿--欧拉递推（RNEA）：\texorpdfstring{$O(n)$}{O(n)} 快算逆动力学}
\label{sec:ad-rnea}

\paragraph{反面：Lagrange 法算起来慢。}
拉格朗日法\emph{看结构}无敌，但\emph{算}起来糟糕：显式组装 $\mathbf{M}$ 已是 $O(n^2)$ 个元素、每个还要累加各杆雅可比；Christoffel 符号有 $O(n^3)$ 个偏导。对实时控制（每周期算一次逆动力学）这不可接受。\textbf{递归牛顿--欧拉算法（RNEA）}给出 $O(n)$ 的逆动力学，且\emph{根本不显式形成 $\mathbf{M},\mathbf{C}$}——它是机器人实时力矩控制的算法基石\cite{featherstone2008rbda}。本节用 P7 已建的\emph{空间向量代数}（\cref{sec:qk-plucker}）来表述它，记号最简洁、对定基/浮动基统一。

\paragraph{借用 P7 的空间向量代数（不重建）。}
RNEA 的现代表述把“$3$ 维力 + $3$ 维力矩”“$3$ 维线速度 + $3$ 维角速度”各塞进\emph{一个}六维 Plücker 空间向量。这套代数本书已在 \cref{sec:qk-plucker} 为四足整机递推建好，本章\emph{直接借用}（编写规范跨章去重：同一记号只建一次）。需要的记号速记：
\begin{itemize}
  \item \textbf{空间速度}（twist）$\mathbf{v}_i\in\mathbb{R}^6$、\textbf{空间加速度} $\mathbf{a}_i\in\mathbb{R}^6$、\textbf{空间力}（wrench）$\mathbf{f}_i\in\mathbb{R}^6$。
  \item \textbf{空间变换} $\mathbf{X}_{i,\lambda(i)}$：把父杆 $\lambda(i)$ 坐标系下的运动向量变到杆 $i$ 坐标系（即 Plücker 变换 $\mathbf{X}$，\cref{eq:qk-plucker-X}）；力向量用其对偶 $\mathbf{X}^*$。
  \item \textbf{关节运动子空间} $\mathbf{S}_i\in\mathbb{R}^{6\times k_i}$：第 $i$ 关节允许的运动方向（转动关节 $k_i=1$，$\mathbf{S}_i$ 是该轴的空间轴），使关节贡献的相对空间速度 $=\mathbf{S}_i\dot{q}_i$。
  \item \textbf{空间叉乘} $\mathbf{v}\times$（运动 $\times$ 运动）与 $\mathbf{v}\times^*$（运动 $\times$ 力，对偶），\textbf{空间惯量} $\mathbf{I}_i\in\mathbb{R}^{6\times6}$（杆 $i$ 的质量/质心/转动惯量打包，\cref{sec:qk-plucker} 已定义）。
\end{itemize}

\paragraph{核心思想：力沿运动链“两遍走”。}
RNEA 把逆动力学拆成两遍沿运动链的递推（\cref{fig:ad-rnea}）：
\begin{enumerate}
  \item \textbf{前向（基座 $\to$ 末端）传播运动学}：从基座出发，逐杆累加得到每杆的空间速度 $\mathbf{v}_i$ 与空间加速度 $\mathbf{a}_i$。父杆的运动 $+$ 本关节加进来的运动 $=$ 本杆的运动。
  \item \textbf{后向（末端 $\to$ 基座）传播力}：在每杆上用牛顿--欧拉（“合外力 = 空间惯量 $\times$ 空间加速度 $+$ 速度积偏置”）算出该杆\emph{所受合空间力} $\mathbf{f}_i$，它等于本杆惯性力加上\emph{子杆传回}的力。把 $\mathbf{f}_i$ 投影到关节轴 $\mathbf{S}_i$ 上即得关节力矩 $\tau_i$。
\end{enumerate}

\begin{figure}[H]
\centering
\begin{tikzpicture}[>=latex, font=\small]
  % links as a chain of boxes
  \foreach \i/\x in {0/0, 1/2.0, 2/4.0, 3/6.0} {
    \node[draw, rounded corners, minimum width=1.3cm, minimum height=0.7cm, fill=main!8] (L\i) at (\x,0) {$\text{link}\,\i$};
  }
  \node[anchor=east] at (L0.west) {基座};
  \node[draw, circle, fill=second!15, inner sep=1pt] at (5.0,0) {};
  % forward arrows (top): velocity/acceleration
  \draw[->, thick, main!70!black] (-0.9,0.95) -- (6.9,0.95)
     node[midway, above]{前向：传播 $\mathbf{v}_i=\mathbf{X}_{i,\lambda}\mathbf{v}_{\lambda}+\mathbf{S}_i\dot{q}_i$，\ $\mathbf{a}_i$（基座 $\to$ 末端）};
  % backward arrows (bottom): force
  \draw[<-, thick, second!70!black] (-0.9,-0.95) -- (6.9,-0.95)
     node[midway, below]{后向：传播 $\mathbf{f}_i=\mathbf{I}_i\mathbf{a}_i+\mathbf{v}_i\times^*\mathbf{I}_i\mathbf{v}_i+\sum_{\text{子}}\mathbf{X}^*\mathbf{f}_j$，\ $\tau_i=\mathbf{S}_i^\top\mathbf{f}_i$（末端 $\to$ 基座）};
  \draw[->] (L0) -- (L1); \draw[->] (L1) -- (L2); \draw[->] (L2) -- (L3);
\end{tikzpicture}
\caption{RNEA 的“两遍递推”：前向沿链传播运动学（速度/加速度），后向沿链回传力（力/力矩）。每杆只与其父、子交互，故总代价 $O(n)$。}
\label{fig:ad-rnea}
\end{figure}

\begin{algo}[递归牛顿--欧拉逆动力学 RNEA]
\label{alg:ad-rnea}
\textbf{输入} $(\mathbf{q},\dot{\mathbf{q}},\ddot{\mathbf{q}})$，各杆空间惯量 $\mathbf{I}_i$、关节子空间 $\mathbf{S}_i$；\textbf{输出} $\boldsymbol{\tau}$。\\
\textbf{重力技巧}：把基座空间加速度初始化为 $\mathbf{a}_0=-\mathbf{a}_g$（$\mathbf{a}_g$ 为重力加速度对应的空间向量），即让“基座向上加速”，则重力\emph{自动}并入加速度递推、无需单列 $\mathbf{g}(\mathbf{q})$\cite{lynch2017modern,featherstone2008rbda}。
\begin{enumerate}
  \item \textbf{前向}（$i=1\to n$，$\lambda(i)$ 为 $i$ 的父）：
  \begin{align}
    \mathbf{v}_i&=\mathbf{X}_{i,\lambda(i)}\,\mathbf{v}_{\lambda(i)}+\mathbf{S}_i\dot{q}_i, \label{eq:ad-rnea-v}\\
    \mathbf{a}_i&=\mathbf{X}_{i,\lambda(i)}\,\mathbf{a}_{\lambda(i)}+\mathbf{S}_i\ddot{q}_i+\mathbf{v}_i\times\mathbf{S}_i\dot{q}_i. \label{eq:ad-rnea-a}
  \end{align}
  （$\mathbf{v}_0=\mathbf{0}$，$\mathbf{a}_0=-\mathbf{a}_g$。\cref{eq:ad-rnea-a} 末项 $\mathbf{v}_i\times\mathbf{S}_i\dot{q}_i$ 是关节速度引起的速度积/科氏效应。）
  \item \textbf{每杆牛顿--欧拉}（合空间力 = 惯性力 + 速度积偏置）：
  \begin{equation}
    \mathbf{f}_i^{\,\mathrm{net}}=\mathbf{I}_i\,\mathbf{a}_i+\mathbf{v}_i\times^{*}\mathbf{I}_i\,\mathbf{v}_i. \label{eq:ad-rnea-fnet}
  \end{equation}
  （$\mathbf{v}\times^*\mathbf{I}\mathbf{v}$ 是“离心/陀螺”偏置力，对应 \cref{eq:ad-coriolis-matrix} 的 $\mathbf{C}\dot{\mathbf{q}}$ 在空间向量下的化身。）
  \item \textbf{后向}（$i=n\to1$）回传子杆力并取关节力矩：
  \begin{align}
    \mathbf{f}_i&=\mathbf{f}_i^{\,\mathrm{net}}+\!\!\sum_{j\in\text{children}(i)}\!\!\mathbf{X}^{*}_{j,i}\,\mathbf{f}_j\;(-\,\mathbf{f}_i^{\mathrm{ext}}), \label{eq:ad-rnea-f}\\
    \tau_i&=\mathbf{S}_i^\top\,\mathbf{f}_i. \label{eq:ad-rnea-tau}
  \end{align}
  （串联臂每杆至多一个子，求和退化为单项。$\mathbf{f}_i^{\mathrm{ext}}$ 为施于杆 $i$ 的外接触力，若有则减去——这是接触/力控的入口。）
\end{enumerate}
每步只做常数次 $6\times6$ 矩阵--向量运算、且只与父/子交互，故\textbf{总复杂度 $O(n)$}。
\end{algo}

\begin{insight}[一次 RNEA = 整条逆动力学 $\mathbf{M}\ddot{\mathbf{q}}+\mathbf{C}\dot{\mathbf{q}}+\mathbf{g}$]
\label{ins:ad-rnea-onepass}
\cref{alg:ad-rnea} 输出的 $\boldsymbol{\tau}$ 正是 \cref{eq:ad-manip} 的左边\emph{整体}——它\emph{同时}算出了 $\mathbf{M}\ddot{\mathbf{q}}$、$\mathbf{C}\dot{\mathbf{q}}$、$\mathbf{g}$ 三项之和，却\emph{不曾显式形成}其中任何一个矩阵。这就是计算力矩控制的工程命脉：把“期望加速度 + PD 修正”当作 $\ddot{\mathbf{q}}$ 喂进 RNEA，一次递推就得到该发的力矩 $\boldsymbol{\tau}=\mathrm{rnea}(\mathbf{q},\dot{\mathbf{q}},\ddot{\mathbf{q}}_d+\mathbf{K}_p\mathbf{e}+\mathbf{K}_d\dot{\mathbf{e}})$（\cref{ch:arm-ctrl} 与 \cref{ch:arm-prac} 详述）。需要单独的 $\mathbf{g}(\mathbf{q})$？令 $\dot{\mathbf{q}}=\ddot{\mathbf{q}}=\mathbf{0}$ 跑一遍 RNEA 即得（这就是 \texttt{computeGeneralizedGravity} 的本质）。需要 $\mathbf{C}\dot{\mathbf{q}}+\mathbf{g}$（非线性偏置）？令 $\ddot{\mathbf{q}}=\mathbf{0}$ 即得。
\end{insight}

\paragraph{与 Lagrange 法的等价。}
两条路推的是\emph{同一个}方程——这不是巧合，而是力学的两种等价表述（牛顿力学与解析力学）的必然。

\begin{insight}[殊途同归：能量观与递推观的分工]
\label{ins:ad-equiv}
\textbf{Euler--Lagrange 与牛顿--欧拉给出完全相同的运动方程 \cref{eq:ad-manip}}，分工互补：
\begin{itemize}
  \item \emph{Euler--Lagrange}（能量观）——内部约束力自动消去，方程结构清晰，最适合\emph{看清结构、证性质}（$\mathbf{M}\succ0$、斜对称、线性可参数化的证明都最自然）。代价：显式组装慢（$O(n^2)\sim O(n^3)$）。
  \item \emph{牛顿--欧拉/RNEA}（递推观）——沿链两遍递推，最适合\emph{快速计算}（$O(n)$ 逆动力学，不形成矩阵）。代价：递推过程不直接“看见”整体结构。
\end{itemize}
工程实践：\textbf{用 Lagrange 法（或本章的结构性质）理解、用 RNEA（Pinocchio \texttt{rnea}）计算}。Modern Robotics \cite{lynch2017modern} 用螺旋/伴随表述 RNEA、Featherstone \cite{featherstone2008rbda} 用 Plücker 空间向量表述，二者是同一算法的两种记号（本书取后者以接续 P7）。
\end{insight}

\begin{note}[螺旋/伴随表述与 Plücker 表述是同一算法]
\label{note:ad-mr-vs-fw}
Modern Robotics（Lynch--Park）把 RNEA 写成：前向 $\mathcal{V}_i=\mathrm{Ad}_{T_{i,i-1}}\mathcal{V}_{i-1}+\mathcal{A}_i\dot{\theta}_i$、$\dot{\mathcal{V}}_i=\mathrm{Ad}_{T_{i,i-1}}\dot{\mathcal{V}}_{i-1}+[\mathrm{ad}_{\mathcal{V}_i}]\mathcal{A}_i\dot{\theta}_i+\mathcal{A}_i\ddot{\theta}_i$；后向 $\mathcal{F}_i=\mathrm{Ad}^\top_{T_{i+1,i}}\mathcal{F}_{i+1}+\mathcal{G}_i\dot{\mathcal{V}}_i-[\mathrm{ad}_{\mathcal{V}_i}]^\top\mathcal{G}_i\mathcal{V}_i$、$\tau_i=\mathcal{F}_i^\top\mathcal{A}_i$\cite{lynch2017modern}。把 $\mathcal{A}_i\!\to\!\mathbf{S}_i$、$\mathrm{Ad}\!\to\!\mathbf{X}$、$[\mathrm{ad}]\!\to\!(\times)$、$\mathcal{G}\!\to\!\mathbf{I}$、对偶变换 $\mathrm{Ad}^\top\!\to\!\mathbf{X}^*$，即逐字对上 \cref{alg:ad-rnea}。本书统一用 P7 的 Plücker 记号（\cref{sec:qk-plucker}）以保持全书一致，但读 Modern Robotics 时按本表对照即可。两套伴随/Plücker 分块顺序的差异见 \cref{note:qk-plucker-order}。
\end{note}

% ════════════════════════════════════════════════════════════════
\section{质量阵与复合刚体算法（CRBA）}
\label{sec:ad-crba}

\paragraph{动机：什么时候非得显式要 $\mathbf{M}$ 不可。}
RNEA 算\emph{逆}动力学不需要显式 $\mathbf{M}$，那为什么还要单独算质量阵？因为有些场合\emph{绕不过} $\mathbf{M}$ 本身：(i) \textbf{正动力学/仿真}要解 $\ddot{\mathbf{q}}=\mathbf{M}^{-1}(\cdots)$；(ii) \textbf{操作空间动力学}（\cref{sec:ad-osd}）要 $\boldsymbol{\Lambda}=(\mathbf{J}\mathbf{M}^{-1}\mathbf{J}^\top)^{-1}$、动力学一致逆要 $\mathbf{M}^{-1}\mathbf{J}^\top$；(iii) \textbf{条件数分析}（\cref{sec:ad-conditioning}）要看 $\mathbf{M}$ 的特征值。\textbf{复合刚体算法（CRBA）}以 $O(n^2)$ 显式算出对称正定的 $\mathbf{M}(\mathbf{q})$，是 Pinocchio \texttt{crba}、RBDL 的标准实现\cite{featherstone2008rbda}。

\paragraph{核心思想：“复合刚体”惯量。}
$\mathbf{M}$ 的第 $(i,j)$ 元 $M_{ij}=\partial^2 K/\partial\dot{q}_i\partial\dot{q}_j$ 度量“关节 $i$ 与关节 $j$ 的惯性耦合”。CRBA 的妙处：把从关节 $i$ 往末端的\emph{所有杆}视为一个\emph{刚性焊死}的“复合刚体”，其空间惯量 $\mathbf{I}^c_i$ 由“子树惯量逐级回传相加”得到；$M_{ij}$ 则由“复合刚体惯量沿关节轴投影 + 沿链上传”读出。直觉：当只有关节 $i$ 单独运动（其余冻结），从 $i$ 往外整段是一块刚体，其对关节 $i$ 的有效惯量就是 $\mathbf{S}_i^\top\mathbf{I}^c_i\mathbf{S}_i$。

\begin{algo}[复合刚体算法 CRBA（算 $\mathbf{M}(\mathbf{q})$）]
\label{alg:ad-crba}
\textbf{输入} $\mathbf{q}$、各杆空间惯量 $\mathbf{I}_i$、子空间 $\mathbf{S}_i$；\textbf{输出} 对称正定 $\mathbf{M}\in\mathbb{R}^{n\times n}$。
\begin{enumerate}
  \item \textbf{复合惯量回传}（$i=n\to1$）：初始化 $\mathbf{I}^c_i=\mathbf{I}_i$，向父累加
  \begin{equation}
    \mathbf{I}^c_{\lambda(i)}\mathrel{+}=\mathbf{X}^{*}_{i,\lambda(i)}\,\mathbf{I}^c_i\,\mathbf{X}_{i,\lambda(i)}. \label{eq:ad-crba-Ic}
  \end{equation}
  \item \textbf{读取矩阵元}（$i=1\to n$）：先算“惯性力” $\mathbf{F}=\mathbf{I}^c_i\mathbf{S}_i$，对角块
  \begin{equation}
    M_{ii}=\mathbf{S}_i^\top\mathbf{F}; \label{eq:ad-crba-diag}
  \end{equation}
  再沿链向根上传 $\mathbf{F}$，对每个祖先 $j$（$j=\lambda(i),\lambda(\lambda(i)),\dots$）：
  \begin{equation}
    \mathbf{F}\leftarrow\mathbf{X}^{*}_{j+,\,j}\,\mathbf{F},\qquad M_{ij}=\mathbf{F}^\top\mathbf{S}_j,\qquad M_{ji}=M_{ij}. \label{eq:ad-crba-offdiag}
  \end{equation}
  （$\mathbf{X}^*_{j+,j}$ 表沿链把力变到祖先 $j$ 的坐标系。串联臂的祖先链就是一条直线。）
\end{enumerate}
对称性 $M_{ji}=M_{ij}$ 由算法\emph{显式}保证（只算上三角、镜像下三角），与 \cref{ins:ad-M-spd} 的解析对称一致。复杂度 $O(n^2)$（每个 $i$ 沿链上传至多 $n$ 步）。
\end{algo}

\begin{insight}[CRBA 的对称正定 = Euler--Lagrange 结论的算法级再现]
\label{ins:ad-crba-struct}
\cref{sec:ad-lagrange} 从动能二次型\emph{解析地}证了 $\mathbf{M}$ 对称正定（\cref{ins:ad-M-spd}）；CRBA \emph{算法地}再现了它：对称由 \cref{eq:ad-crba-offdiag} 的镜像直接保证；正定由“复合刚体空间惯量 $\mathbf{I}^c_i\succ0$ 沿轴投影 $\mathbf{S}_i^\top\mathbf{I}^c_i\mathbf{S}_i>0$”保证（每个对角块为正、且整体是合法动能的 Gram 阵）。\textbf{这正是 P4 \cref{prop:ctrl-robot-props} 性质 1 的“算法来历”}：P4 把 $\mathbf{M}\succ0$ 当已知前提用，本节给出它怎么\emph{算}、为什么\emph{算出来必然}对称正定。配合 \cref{prop:ad-skew}（斜对称），P4 的两条承重性质在本章都有了完整的“出生证明”。
\end{insight}

\begin{note}[三个算法的家谱：RNEA / CRBA / ABA]
\label{note:ad-algo-family}
本章出现的多体动力学算法可按“解什么问题 + 复杂度”归位\cite{featherstone2008rbda}：
\begin{itemize}
  \item \textbf{RNEA}（\cref{alg:ad-rnea}）：逆动力学 $\boldsymbol{\tau}=\mathrm{ID}(\mathbf{q},\dot{\mathbf{q}},\ddot{\mathbf{q}})$，$O(n)$，不形成矩阵——\emph{控制/规划}主力。
  \item \textbf{CRBA}（\cref{alg:ad-crba}）：显式 $\mathbf{M}(\mathbf{q})$，$O(n^2)$——需要质量阵\emph{本身}时（仿真、$\boldsymbol{\Lambda}$、条件数）用。
  \item \textbf{ABA}（关节体算法，本章仅点名）：正动力学 $\ddot{\mathbf{q}}=\mathrm{FD}(\mathbf{q},\dot{\mathbf{q}},\boldsymbol{\tau})$，$O(n)$，不显式求 $\mathbf{M}^{-1}$——\emph{仿真}主力，详见 \cite{featherstone2008rbda}。
\end{itemize}
工程上常见组合：仿真用 ABA；控制用 RNEA（前馈）+ CRBA（操作空间惯量）。Pinocchio 三者俱全。
\end{note}

% ════════════════════════════════════════════════════════════════
\section{操作空间动力学}
\label{sec:ad-osd}

\paragraph{动机：在“任务空间”里写动力学。}
许多任务的目标天然在\emph{末端任务空间}（笛卡尔位置/姿态），而非关节空间——“末端走直线”“末端顶恒力”“末端呈现某刚度”。上一章给了\emph{运动学}对偶：末端速度 $\dot{\mathbf{x}}=\mathbf{J}\dot{\mathbf{q}}$、静力学对偶 $\boldsymbol{\tau}=\mathbf{J}^\top\mathbf{F}$（\cref{sec:ak-jacobian}）。但运动学对偶\emph{不够}：它告诉你“末端力怎么折成关节力矩”，却没告诉你“要让末端产生期望加速度 $\ddot{\mathbf{x}}$，该施多大末端力 $\mathbf{F}$”——后者要\emph{动力学}。Khatib（1987）的\textbf{操作空间公式}把关节动力学\emph{投影}到任务空间，给出任务空间里的牛顿第二定律\cite{khatib1987operational}。这是后续操作空间控制（\cref{ch:arm-ctrl}）的地基；本节只建\emph{动力学}（控制律留到 \cref{ch:arm-ctrl}）。

\paragraph{推导：把关节动力学投影到任务空间。}

\begin{derivation}[操作空间惯量 $\boldsymbol{\Lambda}=(\mathbf{J}\mathbf{M}^{-1}\mathbf{J}^\top)^{-1}$ 的来历]
从关节动力学 $\mathbf{M}\ddot{\mathbf{q}}+\mathbf{C}\dot{\mathbf{q}}+\mathbf{g}=\boldsymbol{\tau}=\mathbf{J}^\top\mathbf{F}$（末端力 $\mathbf{F}$ 经对偶折成关节力矩）出发。左乘 $\mathbf{J}\mathbf{M}^{-1}$：
\[
  \mathbf{J}\ddot{\mathbf{q}}+\mathbf{J}\mathbf{M}^{-1}(\mathbf{C}\dot{\mathbf{q}}+\mathbf{g})=\mathbf{J}\mathbf{M}^{-1}\mathbf{J}^\top\mathbf{F}.
\]
对 $\dot{\mathbf{x}}=\mathbf{J}\dot{\mathbf{q}}$ 求时间导得 $\ddot{\mathbf{x}}=\mathbf{J}\ddot{\mathbf{q}}+\dot{\mathbf{J}}\dot{\mathbf{q}}\Rightarrow\mathbf{J}\ddot{\mathbf{q}}=\ddot{\mathbf{x}}-\dot{\mathbf{J}}\dot{\mathbf{q}}$。代入：
\[
  \ddot{\mathbf{x}}-\dot{\mathbf{J}}\dot{\mathbf{q}}+\mathbf{J}\mathbf{M}^{-1}(\mathbf{C}\dot{\mathbf{q}}+\mathbf{g})=(\mathbf{J}\mathbf{M}^{-1}\mathbf{J}^\top)\mathbf{F}.
\]
定义\emph{操作空间惯量} $\boldsymbol{\Lambda}\triangleq(\mathbf{J}\mathbf{M}^{-1}\mathbf{J}^\top)^{-1}$（$\mathbf{J}$ 行满秩时存在），两边左乘 $\boldsymbol{\Lambda}$ 并移项：
\[
  \boldsymbol{\Lambda}\ddot{\mathbf{x}}+\underbrace{\boldsymbol{\Lambda}\big(\mathbf{J}\mathbf{M}^{-1}\mathbf{C}\dot{\mathbf{q}}-\dot{\mathbf{J}}\dot{\mathbf{q}}\big)}_{\boldsymbol{\mu}\ \text{任务空间科氏/离心}}+\underbrace{\boldsymbol{\Lambda}\mathbf{J}\mathbf{M}^{-1}\mathbf{g}}_{\mathbf{p}\ \text{任务空间重力}}=\mathbf{F}.
\]
\end{derivation}

\noindent 得到\emph{任务空间动力学}（与 P4 \cref{eq:ctrl-osc} 一致）：
\begin{equation}
  \boxed{\;\boldsymbol{\Lambda}(\mathbf{q})\ddot{\mathbf{x}}+\boldsymbol{\mu}(\mathbf{q},\dot{\mathbf{q}})+\mathbf{p}(\mathbf{q})=\mathbf{F},\qquad \boldsymbol{\Lambda}=(\mathbf{J}\mathbf{M}^{-1}\mathbf{J}^\top)^{-1}.\;}
  \label{eq:ad-osc}
\end{equation}
形式与关节空间 $\mathbf{M}\ddot{\mathbf{q}}+\mathbf{C}\dot{\mathbf{q}}+\mathbf{g}=\boldsymbol{\tau}$ \emph{完全平行}：$\boldsymbol{\Lambda}$ 是“末端感受到的有效惯量”，$\boldsymbol{\mu},\mathbf{p}$ 是任务空间的科氏/重力，$\mathbf{F}$ 是末端广义力。关节力矩仍由 $\boldsymbol{\tau}=\mathbf{J}^\top\mathbf{F}$ 给出。

\begin{insight}[$\boldsymbol{\Lambda}$ 是“末端的等效质量”——为什么需要它]
\label{ins:ad-Lambda-meaning}
为什么不能只用 $\boldsymbol{\tau}=\mathbf{J}^\top\mathbf{F}$、非要 $\boldsymbol{\Lambda}$？因为同一个末端力 $\mathbf{F}$，在不同位形会引起\emph{不同}的末端加速度——末端“感觉到”的惯量随位形和方向变化（手臂伸直时末端沿轴向“重”、横向“轻”）。$\boldsymbol{\Lambda}=(\mathbf{J}\mathbf{M}^{-1}\mathbf{J}^\top)^{-1}$ 精确刻画这个\emph{方向相关、位形相关}的等效惯量。有了它，控制律 $\mathbf{F}=\hat{\boldsymbol{\Lambda}}\mathbf{F}^*+\hat{\boldsymbol{\mu}}+\hat{\mathbf{p}}$ 才能用前馈\emph{抵消}末端的真实惯量，得到解耦的任务空间行为 $\ddot{\mathbf{x}}=\mathbf{F}^*$——这就是操作空间控制（\cref{ch:arm-ctrl} 完整推导）。注意 \cref{eq:ad-osc} 的核心是 $\mathbf{M}^{-1}$：$\boldsymbol{\Lambda}$ 的可算性、数值稳定性\emph{完全系于} $\mathbf{M}$ 的条件数——这道伏笔在下一节引爆。
\end{insight}

\paragraph{动力学一致广义逆与零空间。}
冗余臂（关节数 $>$ 任务维）有“零空间”：不影响末端的内部自运动。如何把“关节力矩”分解成“管任务的”加“管零空间的”而\emph{互不干扰}？这要\emph{动力学一致}的广义逆。

\begin{definition}[动力学一致广义逆 $\bar{\mathbf{J}}$ 与一致零空间投影]
\label{def:ad-dyn-consistent}
\rebuilt{（据 Khatib \cite{khatib1987operational}）}\emph{动力学一致雅可比广义逆}
\begin{equation}
  \bar{\mathbf{J}}\triangleq\mathbf{M}^{-1}\mathbf{J}^\top\boldsymbol{\Lambda}=\mathbf{M}^{-1}\mathbf{J}^\top(\mathbf{J}\mathbf{M}^{-1}\mathbf{J}^\top)^{-1}.
  \label{eq:ad-Jbar}
\end{equation}
对应的\emph{动力学一致零空间投影}
\begin{equation}
  \mathbf{N}=\mathbf{I}-\mathbf{J}^\top\bar{\mathbf{J}}^\top=\mathbf{I}-\mathbf{J}^\top(\mathbf{J}\mathbf{M}^{-1}\mathbf{J}^\top)^{-1}\mathbf{J}\mathbf{M}^{-1}.
  \label{eq:ad-N-dyn}
\end{equation}
其定义性质：$\bar{\mathbf{J}}$ 是\emph{唯一}使“任意零空间关节力矩 $\mathbf{N}\boldsymbol{\tau}_0$ 对末端产生\emph{零加速度}”的广义逆；等价地，它给出\emph{最小瞬时动能}的关节速度解。
\end{definition}

\begin{insight}[动力学一致 vs 运动学（Moore--Penrose）：差在一个 $\mathbf{M}$]
\label{ins:ad-consistent-vs-mp}
上一章用的是\emph{运动学}伪逆 $\mathbf{J}^+=\mathbf{J}^\top(\mathbf{J}\mathbf{J}^\top)^{-1}$（Moore--Penrose）、对应零空间投影 $\mathbf{N}_{\text{kin}}=\mathbf{I}-\mathbf{J}^+\mathbf{J}$（\cref{sec:ak-redundancy}）。本节的\emph{动力学一致}逆 $\bar{\mathbf{J}}=\mathbf{M}^{-1}\mathbf{J}^\top\boldsymbol{\Lambda}$ 多了 $\mathbf{M}^{-1}$ 加权——这\emph{不是}吹毛求疵：用运动学投影做\emph{力矩级}零空间控制时，零空间的避奇异/关节中心化力矩会“\emph{泄漏}”出一个末端扰动加速度，破坏任务纯净；只有动力学一致逆能保证零空间力矩对末端\emph{零加速度}。这正是许多工业实现（图省事用 $\mathbf{J}^+$）阻抗不纯的工程缘由（P4 \cref{sec:ctrl-robot-osc} 与 \cref{ch:force-wbc} 已强调“两种伪逆严格区分”）。\textbf{但代价是引入了 $\mathbf{M}^{-1}$}——当 $\mathbf{M}$ 病态，动力学一致逆反而\emph{更脆弱}（下一节的核心实证）。
\end{insight}

% ════════════════════════════════════════════════════════════════
\section{质量阵条件数与数值稳定性}
\label{sec:ad-conditioning}

\paragraph{动机：从“可逆”到“求逆数值稳定”——一道致命的鸿沟。}
$\mathbf{M}(\mathbf{q})\succ0$（\cref{ins:ad-M-spd}）保证它\emph{数学上处处可逆}。但本章所有依赖 $\mathbf{M}^{-1}$ 的构造——正动力学、操作空间惯量 $\boldsymbol{\Lambda}=(\mathbf{J}\mathbf{M}^{-1}\mathbf{J}^\top)^{-1}$、动力学一致逆 $\bar{\mathbf{J}}=\mathbf{M}^{-1}\mathbf{J}^\top\boldsymbol{\Lambda}$——在\emph{数值上}稳不稳，取决于 $\mathbf{M}$ 离奇异有多近，由\emph{条件数}度量。本节用第一性原理（条件数、自然时间常数）把“病态质量阵”讲透，并以一台真实异质惯量臂的穷尽实证（源自 \cref{ch:arm-prac} 的工程根因分析）坐实它。\textbf{这是本部区别于纯教科书的最深实证之一}。

\begin{definition}[质量阵条件数]
\label{def:ad-condM}
对称正定 $\mathbf{M}$ 的\emph{条件数}（$2$-范数）$\mathrm{cond}(\mathbf{M})=\lambda_{\max}(\mathbf{M})/\lambda_{\min}(\mathbf{M})$（最大/最小特征值之比）。$\mathrm{cond}=1$ 最良态（各方向惯量均衡）；$\mathrm{cond}\gg1$ \emph{病态}——$\mathbf{M}$ 在某方向“极重”、另一方向“极轻”，求 $\mathbf{M}^{-1}$ 时轻方向的微小数值误差被放大 $\mathrm{cond}$ 倍。
\end{definition}

\paragraph{第一性原理：异质惯量臂为何\emph{固有}病态。}
质量阵的对角元 $M_{ii}=\mathbf{S}_i^\top\mathbf{I}^c_i\mathbf{S}_i$ 是“关节 $i$ 的有效惯量”（\cref{eq:ad-crba-diag}）。若机械臂\emph{重肩肘 + 轻腕}——肩肘 $M_{22},M_{33}$ 大（驱动远端整条臂，等效惯量大），腕 $M_{44},M_{55},M_{66}$ 小（只驱动末端小法兰，真实 CAD 惯量量级 $\sim2\times10^{-4}$）——则 $\mathbf{M}$ 的特征值\emph{天然}横跨好几个数量级，$\mathrm{cond}(\mathbf{M})$ 必然大。\textbf{这不是建模 bug，是“重肩肘 + 轻腕”这一\emph{本体设计}的固有属性}：条件数由本体惯量结构决定，换参数（正则化）治标不治本，换层（速度级 vs 力矩级）才治本。

\begin{insight}[实证 1：病态质量阵让力矩级 OSC 零空间收不住末端]
\label{ins:ad-osc-illcond}
一台六轴力矩臂（\cref{ch:arm-prac} 的 \texttt{arm\_dm}，腕惯量 $\sim2\times10^{-4}$）上的穷尽验证给出一条触目惊心的结论。任务：用动力学一致零空间投影 \cref{eq:ad-N-dyn} 让末端\emph{锁死}、同时让手肘在零空间自运动。结果：自运动出现了（J2/J3 重构 $1.1$--$1.4\,\mathrm{rad}$），但\textbf{末端没收住（$\|\Delta\mathbf{x}\|\approx70$--$110\,\mathrm{mm}$，而非 $<10\,\mathrm{mm}$）}。四次定位尝试逐一排除：(1) 运动学投影 $\mathbf{N}=\mathbf{I}-\mathbf{J}^+\mathbf{J}$ 在力矩级泄漏 $85\,\mathrm{mm}$；(2) 精确动力学投影 \cref{eq:ad-N-dyn} \emph{数值爆炸}（腕 $2\times10^{-4}$ 惯量使 $\mathbf{M}^{-1}$ 巨大、奇异处 $(\mathbf{J}\mathbf{M}^{-1}\mathbf{J}^\top)^{-1}$ 爆，$\|\boldsymbol{\tau}_{\text{null}}\|\to18000$）；(3) 正则化 $\mathbf{M}+\lambda\mathbf{I}$ / 阻尼逆 → 稳但末端仍漂 $70\,\mathrm{mm}$；(4) 锁腕 → $107\,\mathrm{mm}$，均无改善。\textbf{真根因 = 质量阵病态}：实测 $\mathrm{cond}(\mathbf{M})\approx527$--$1625$（随位形变）。精确投影数值不稳；为稳定加的正则化\emph{破坏了 $\mathbf{N}$ 的投影性}——投影阵特征值本应为 $\{0,1\}$，正则后不再是，于是放大零空间力矩 $\boldsymbol{\tau}_0$ 漏入任务空间（$10$--$50\times$）→ 末端漂。
\end{insight}

\begin{insight}[实证 2：占位惯量是“条件数拐杖”，真几何更病态]
\label{ins:ad-placeholder}
更深一层：上述腕端 link6 的\emph{真实几何}惯量比仿真用的占位值\emph{更小}。trimesh 真几何分析（\cref{ch:arm-prac}）显示 link6 是 $1$--$2\,\mathrm{cm}$ 小法兰（体积 $\sim3\times10^{-6}\,\mathrm{m}^3$），真几何惯量 $\sim5\times10^{-6}$——比仿真里填的 $1\times10^{-3}$ 占位值\emph{小 $200\times$}。后果：裸真几何下 $\mathrm{cond}(\mathbf{M})$ \emph{爆到 $1.3\times10^4$--$5.6\times10^4$}。\textbf{那个 $1\times10^{-3}$ 占位惯量，实质是一根“条件数拐杖”}——它（连同 link6 质量从 $0.003$ 撑到 $0.15\,\mathrm{kg}$）静默地把 $\mathrm{cond}$ 压到 $\sim10^3$ 量级，让动力学算得动。即便换一个好条件数的真末端（如 $0.3\,\mathrm{kg}$ gripper 把 $\mathrm{cond}$ 砍半到 $300+$），OSC 仍漂 $62\,\mathrm{mm}$ 且自运动塌缩——而 link4/link5 真 CAD 惯量 $2$--$3\times10^{-4}$ 是\emph{固有地板}、不根治。教训：\emph{用真几何惯量做条件数体检，别被“能跑”的占位值蒙蔽}。
\end{insight}

\paragraph{第二条战线：低惯量关节与控制频率——自然时间常数。}
病态质量阵伤的是\emph{空间}（求逆方向放大误差）；低惯量关节还伤\emph{时间}——它的“自然时间常数”可能快过控制频率，导致欠采样发散。

\begin{insight}[自然频率 $\omega_n=\sqrt{k_p/I}$ 判稳：轻腕为何\emph{更难}控]
\label{ins:ad-wn}
一个二阶受控关节（刚度 $k_p$、惯量 $I$）的自然角频率 $\omega_n=\sqrt{k_p/I}$。\textbf{惯量 $I$ 越小，$\omega_n$ 越高}——这反直觉：轻的东西不是“更好推”吗？是更好\emph{推}，但更难\emph{稳定控制}。原因：离散控制要稳定，须采样足够密 $\omega_n\,\Delta t\lesssim2$（每个振荡周期取够多样本）；惯量 $2\times10^{-4}$ 的腕，自然时间常数 $\approx I/\text{damping}\approx2\times10^{-4}/0.1=2\,\mathrm{ms}$（约 $500\,\mathrm{Hz}$），\emph{快于}典型控制率（$100$--$200\,\mathrm{Hz}$）→ \textbf{欠采样、闭环发散}。实证（\cref{ch:arm-prac}）：均匀增益 PD 控全关节会失稳（腕 $\omega_n=\sqrt{150/2\times10^{-4}}\approx866\,\mathrm{rad/s}\gg$ 控制率）；对 $2\times10^{-4}$ 惯量腕施原始 PD，力矩在饱和前竟冲到 $888\,\mathrm{N\!\cdot\! m}$、即便钳到 $26$ 仍产生 $1.3\times10^5\,\mathrm{rad/s^2}$ 加速度——“小惯量末端在控制率下原始 PD 必炸”最有说服力的实证。
\end{insight}

\begin{insight}[第一性原理选型：换层 > 换参，加权 $\mathbf{M}$ 给正确力矩权限]
\label{ins:ad-selection-law}
两条战线汇成一条\emph{选型律}（直接喂给 \cref{ch:arm-ctrl} 的控制率选择）：
\begin{enumerate}
  \item \textbf{异质惯量臂（病态 $\mathbf{M}$）的冗余解析正解是\emph{运动学/速度级}}（\cref{sec:ak-redundancy} 的 $\mathbf{N}=\mathbf{I}-\mathbf{J}^+\mathbf{J}$，\emph{无} $\mathbf{M}^{-1}$，实测末端锁 $0.00\,\mathrm{mm}$），\emph{力矩级 OSC（\cref{eq:ad-N-dyn}）只在良态臂可行}。别在力矩层硬刚一个本体注定病态的算法——\emph{换层比换参数（正则化）更对}。
  \item \textbf{低惯量腕必须解耦}：大范围计算力矩跟踪时把腕的命令加速度置零、另行温和保持（靠物理阻尼兜底），否则腕误差经\emph{质量阵耦合项} $M_{ij}$ 注入主关节力矩、腐蚀肩肘跟踪（\cref{ch:arm-prac} 实证：去耦合后 $\tau_2$ 从异常的 $15$ 回到正常的 $3.5$）。
  \item \textbf{计算力矩的 $\mathbf{M}(\mathbf{q})$ 加权天然给每关节正确力矩权限}——这是异质惯量臂\emph{用逆动力学控制}（而非均匀增益 PD）的第二理由：$\mathbf{M}$ 自动按有效惯量分配力矩，重肩肘得大力矩、轻腕得小力矩。
\end{enumerate}
一句话：\textbf{用第一性原理（条件数、自然时间常数）判断算法可行性，而非在出问题的层级反复调参}。
\end{insight}

\begin{pitfall}[概念误区：“正定即可逆即可放心求逆”]
\label{pit:ad-spd-not-safe}
新手见 $\mathbf{M}\succ0$ 便认定“处处可逆，求逆无忧”。错。\textbf{“数学可逆”与“数值稳定求逆”是两件事}——$\mathbf{M}$ 永不奇异（$\lambda_{\min}>0$），但 $\lambda_{\min}$ 可以小到 $10^{-4}$，使 $\mathrm{cond}(\mathbf{M})\sim10^3$--$10^4$，求 $\mathbf{M}^{-1}$ 时轻方向误差被放大千万倍。这就是为什么照搬冗余臂教科书的 OSC + 动力学一致零空间，在异质惯量臂上\emph{收不住末端}（\cref{ins:ad-osc-illcond}）。\emph{修法不是“调正则化系数”}（那会破坏投影性、保证失败），而是\emph{换控制层}（用速度级冗余解析）或\emph{换本体}（良态惯量）。诊断工具：算 $\mathrm{cond}(\mathbf{M})$（用真几何惯量），$>10^3$ 就要警惕力矩级算法。
\end{pitfall}

% ════════════════════════════════════════════════════════════════
\section{动力学计算工具速查}
\label{sec:ad-practice}

\paragraph{从理论到代码：Pinocchio 的用途映射。}
本章的算法在工程中无须手写——Pinocchio（C++/Python 多体动力学库，\cref{ch:arm-prac} 全程用它）把 RNEA/CRBA/重力/雅可比都封装好。要点是\emph{把每个 API 对回本章的哪条公式}，知道它\emph{算什么、何时用}，而非记签名。

\begin{codebox}[Python]{Pinocchio 动力学 API <-> 本章公式（用途映射，要点伪代码）}
import pinocchio as pin
# model, data = pin.buildModelFromUrdf(...) ; q, v, a 为关节角/速度/加速度

# 逆动力学 RNEA:tau = M(q) a + C(q,v) v + g(q)  —— 本章 alg:ad-rnea、ins:ad-rnea-onepass
# 计算力矩控制就用它:a_cmd = a_d + Kp e + Kd e_dot，一次得该发力矩
tau = pin.rnea(model, data, q, v, a)              # 一次递推 = 整条逆动力学，O(n)

# 重力力矩 g(q):令 v=a=0 的 RNEA —— 本章 eq:ad-gravity;重力补偿控制 tau = g(q)
g = pin.computeGeneralizedGravity(model, data, q)

# 质量阵 M(q):CRBA —— 本章 alg:ad-crba（对称正定，仿真/Lambda/条件数才需显式 M）
M = pin.crba(model, data, q)                      # O(n^2)，data.M 上三角，需对称化

# 非线性偏置 C(q,v) v + g(q):令 a=0 的 RNEA
nle = pin.nonLinearEffects(model, data, q, v)

# 几何雅可比 J（末端速度 x_dot = J q_dot）—— 上一章 sec:ak-jacobian
J = pin.computeFrameJacobian(model, data, q, ee_id, pin.LOCAL_WORLD_ALIGNED)

# 操作空间惯量 Lambda = (J M^-1 J^T)^-1 —— 本章 eq:ad-osc（病态 M 时此处数值不稳，见 sec:ad-conditioning）
Minv = np.linalg.inv(M); Lambda = np.linalg.inv(J @ Minv @ J.T)
# 条件数体检（务必用真几何惯量）:cond(M) > 1e3 警惕力矩级算法
cond_M = np.linalg.cond(M)
\end{codebox}

\noindent 完整的清洗模型、力矩控制节点、计算力矩 vs PD+重力补偿对比、OSC 病态实证，都在本部实践章 \cref{ch:arm-prac} 一条连续会话里落地。本章只给“理论 $\leftrightarrow$ 工具”的桥。

% ════════════════════════════════════════════════════════════════
\section{常见误解汇总}
\label{sec:ad-misconceptions}

\begin{enumerate}
  \item \textbf{“有了运动学就够了，动力学是锦上添花。”} 错。运动学规划出的轨迹\emph{未必力矩可执行}（\cref{ins:ad-kin-not-enough}）；接触/力控、高速跟踪根本离不开动力学。把“限加速度盒”当力矩约束的代理是危险近似（有效惯量随位形变）。
  \item \textbf{“Euler--Lagrange 和牛顿--欧拉是两种不同的动力学。”} 错。它们是同一力学的两种等价表述，给出\emph{完全相同}的运动方程 \cref{eq:ad-manip}（\cref{ins:ad-equiv}）。分工：Lagrange 看结构、RNEA 快算。
  \item \textbf{“逆动力学要先组装 $\mathbf{M},\mathbf{C},\mathbf{g}$ 再相乘。”} 错且慢。RNEA 一次 $O(n)$ 递推就得到 $\mathbf{M}\ddot{\mathbf{q}}+\mathbf{C}\dot{\mathbf{q}}+\mathbf{g}$ 整体，\emph{不形成}任何矩阵（\cref{ins:ad-rnea-onepass}）。Christoffel 符号是\emph{看结构}用的，不是\emph{算}用的。
  \item \textbf{“科氏阵 $\mathbf{C}$ 是唯一确定的。”} 错。$\mathbf{C}$ 不唯一；斜对称性质（\cref{prop:ad-skew}）只对 Christoffel 选法成立（\cref{pit:ad-coriolis-unique}）。
  \item \textbf{“$\mathbf{M}\succ0$ 正定，所以求 $\mathbf{M}^{-1}$ 总是安全的。”} 错。正定保证\emph{数学}可逆，不保证\emph{数值}稳定；病态 $\mathbf{M}$（$\mathrm{cond}\sim10^3$--$10^4$）会让 OSC/动力学一致逆失效（\cref{pit:ad-spd-not-safe}、\cref{ins:ad-osc-illcond}）。
  \item \textbf{“低惯量的腕关节更好控制（轻嘛）。”} 错。惯量越小自然频率 $\omega_n=\sqrt{k_p/I}$ 越高，越易在控制率下欠采样发散（\cref{ins:ad-wn}）——更难稳控，不是更好控。
  \item \textbf{“动力学一致逆只是 Moore--Penrose 伪逆换个写法。”} 错。二者差一个 $\mathbf{M}^{-1}$ 加权；只有动力学一致逆保证零空间力矩对末端零加速度（\cref{ins:ad-consistent-vs-mp}）。但它引入 $\mathbf{M}^{-1}$，病态时更脆。
\end{enumerate}

% ════════════════════════════════════════════════════════════════
\section{小结}
\label{sec:ad-summary}

本章从两条互补的路（能量观 Euler--Lagrange、递推观牛顿--欧拉 RNEA）推出同一个机械臂运动方程 $\mathbf{M}\ddot{\mathbf{q}}+\mathbf{C}\dot{\mathbf{q}}+\mathbf{g}=\boldsymbol{\tau}$，给出了 P4 \cref{eq:ctrl-manip} 与结构性质 \cref{prop:ctrl-robot-props} 的\emph{完整推导来历}；用 CRBA 算质量阵、建立操作空间动力学，并以第一性原理（条件数、自然时间常数）把“病态质量阵”讲透——后者是本部独有的最深实证。

\paragraph{符号表。}
\begin{center}
\begin{tabular}{ll}
\toprule
符号 & 含义 \\
\midrule
$\mathbf{q},\dot{\mathbf{q}},\ddot{\mathbf{q}}\in\mathbb{R}^n$ & 关节角 / 速度 / 加速度 \\
$\boldsymbol{\tau}\in\mathbb{R}^n$ & 关节驱动力矩 \\
$\mathbf{M}(\mathbf{q})\succ0$ & 质量（惯性）阵，对称正定（\cref{eq:ad-kinetic}）\\
$\mathbf{C}(\mathbf{q},\dot{\mathbf{q}})$ & 科氏/离心阵，Christoffel 选法（\cref{eq:ad-coriolis-matrix}）\\
$\mathbf{g}(\mathbf{q})$ & 重力力矩，$=\partial P/\partial\mathbf{q}$（\cref{eq:ad-gravity}）\\
$c_{ijk}$ & 第一类 Christoffel 符号（\cref{eq:ad-christoffel}）\\
$\mathbf{v}_i,\mathbf{a}_i,\mathbf{f}_i\in\mathbb{R}^6$ & 杆 $i$ 的空间速度/加速度/力（Plücker，\cref{sec:qk-plucker}）\\
$\mathbf{X}_{i,\lambda(i)},\ \mathbf{X}^*$ & 空间变换（运动/力对偶）\\
$\mathbf{S}_i,\ \mathbf{I}_i$ & 关节运动子空间 / 杆空间惯量 \\
$\mathbf{I}^c_i$ & 复合刚体空间惯量（\cref{eq:ad-crba-Ic}）\\
$\mathbf{J},\ \mathbf{J}^+$ & 几何雅可比 / Moore--Penrose 伪逆 \\
$\boldsymbol{\Lambda}=(\mathbf{J}\mathbf{M}^{-1}\mathbf{J}^\top)^{-1}$ & 操作空间惯量（\cref{eq:ad-osc}）\\
$\bar{\mathbf{J}}=\mathbf{M}^{-1}\mathbf{J}^\top\boldsymbol{\Lambda}$ & 动力学一致广义逆（\cref{eq:ad-Jbar}）\\
$\mathbf{N}=\mathbf{I}-\mathbf{J}^\top\bar{\mathbf{J}}^\top$ & 动力学一致零空间投影（\cref{eq:ad-N-dyn}）\\
$\mathrm{cond}(\mathbf{M})=\lambda_{\max}/\lambda_{\min}$ & 质量阵条件数（\cref{def:ad-condM}）\\
$\omega_n=\sqrt{k_p/I}$ & 受控关节自然角频率（\cref{ins:ad-wn}）\\
\bottomrule
\end{tabular}
\end{center}

\paragraph{定理与算法速查。}
\begin{center}
\begin{tabular}{lll}
\toprule
名称 & 内容 & 出处 \\
\midrule
运动方程 \cref{eq:ad-manip} & $\mathbf{M}\ddot{\mathbf{q}}+\mathbf{C}\dot{\mathbf{q}}+\mathbf{g}=\boldsymbol{\tau}$（= P4 \cref{eq:ctrl-manip}）& 本章给推导来历 \\
Euler--Lagrange \cref{eq:ad-el} & $\frac{\mathrm{d}}{\mathrm{d}t}\frac{\partial L}{\partial\dot{q}_i}-\frac{\partial L}{\partial q_i}=\tau_i$ & \cref{sec:ad-lagrange} \\
斜对称 \cref{prop:ad-skew} & $\dot{\mathbf{M}}-2\mathbf{C}$ 斜对称（无源性根）& = \cref{prop:ctrl-robot-props} 性质 2 \\
RNEA \cref{alg:ad-rnea} & $O(n)$ 逆动力学，两遍递推 & \cref{sec:ad-rnea} \\
CRBA \cref{alg:ad-crba} & $O(n^2)$ 算 $\mathbf{M}(\mathbf{q})$，对称正定 & \cref{sec:ad-crba} \\
操作空间动力学 \cref{eq:ad-osc} & $\boldsymbol{\Lambda}\ddot{\mathbf{x}}+\boldsymbol{\mu}+\mathbf{p}=\mathbf{F}$ & = P4 \cref{eq:ctrl-osc} \\
选型律 \cref{ins:ad-selection-law} & 病态臂用速度级、轻腕解耦、$\mathbf{M}$ 加权 & \cref{sec:ad-conditioning} \\
\bottomrule
\end{tabular}
\end{center}

% ════════════════════════════════════════════════════════════════
\section{延伸阅读}
\label{sec:ad-further}
\begin{itemize}
  \item \textbf{Euler--Lagrange / 结构性质}：Spong, Hutchinson, Vidyasagar, \emph{Robot Modeling and Control}\cite{spong2006robot}——动力学章对 Christoffel、无源性、线性可参数化讲得最系统；Tedrake, \emph{Underactuated Robotics}\cite{tedrake_underactuated} 给现代视角与控制衔接。
  \item \textbf{螺旋/伴随表述的牛顿--欧拉}：Lynch, Park, \emph{Modern Robotics}\cite{lynch2017modern} ch.8——用旋量/伴随表述 RNEA，与本书 P0 李群主线最贴合（对照 \cref{note:ad-mr-vs-fw}）。
  \item \textbf{空间向量代数与多体算法（RNEA/CRBA/ABA）}：Featherstone, \emph{Rigid Body Dynamics Algorithms}\cite{featherstone2008rbda}——本章 RNEA/CRBA 的算法源头，含 $O(n)$ 复杂度证明、各算法伪代码逐行；本书 P7 \cref{sec:qk-plucker} 已据它建好空间向量记号。
  \item \textbf{操作空间公式}：Khatib, \emph{A Unified Approach for Motion and Force Control}\cite{khatib1987operational}——操作空间惯量与动力学一致逆的奠基论文。
\end{itemize}

% ════════════════════════════════════════════════════════════════
\section{后续关系}
\label{sec:ad-relations}
\begin{itemize}
  \item \textbf{承上}：本章动能/雅可比来自 \cref{ch:arm-kin}（质心雅可比 $\mathbf{J}_{v_i},\mathbf{J}_{\omega_i}$）、\cref{sec:ak-jacobian}（几何雅可比、静力学对偶）；空间向量代数来自 P7 \cref{sec:qk-plucker}；位姿/伴随来自 P0 \cref{ch:rigid_body,ch:lie}。
  \item \textbf{深化 P4}：本章是 P4 \cref{sec:ctrl-robot-dyn} 的“完整推导来历”——P4 给方程 \cref{eq:ctrl-manip} 与性质 \cref{prop:ctrl-robot-props} 作控制起点，本章证它们怎么来、怎么算（\cref{ins:ad-passivity,ins:ad-crba-struct}）。
  \item \textbf{启下 · 轨迹}：\cref{ch:arm-traj} 的力矩约束时间参数化用本章逆动力学（\cref{alg:ad-rnea}）把 $\boldsymbol{\tau}=\mathbf{M}\ddot{\mathbf{q}}+\mathbf{C}\dot{\mathbf{q}}+\mathbf{g}$ 塞进约束（\cref{ins:ad-kin-not-enough} 的落地）。
  \item \textbf{启下 · 控制}：\cref{ch:arm-ctrl} 的计算力矩用 \cref{ins:ad-rnea-onepass}（一次 RNEA）、操作空间控制 OSC 用本章 \cref{eq:ad-osc} 的任务空间动力学完整推导、异质惯量臂控制率选择用本章 \cref{ins:ad-selection-law} 的条件数选型律。P4 \cref{sec:ctrl-robot-osc} 给良态前提下的理想公式，本章 \cref{sec:ad-conditioning} 给病态本体上的失效边界，二者互补。
  \item \textbf{实践}：\cref{ch:arm-prac} 把本章动力学落到真臂 \texttt{arm\_dm}（Pinocchio \texttt{rnea}/\allowbreak\texttt{crba}、条件数根因分析、计算力矩 vs PD+g）。
  \item \textbf{启下 · RL}：P8 \cref{ch:rlmc-manip}（强化学习操作）散提的质量阵/重力补偿/动力学一致逆，本章提供经典地基。
\end{itemize}

% ════════════════════════════════════════════════════════════════
\section{故障排查}
\label{sec:ad-troubleshooting}

\begin{center}
\begin{tabular}{p{0.30\textwidth} p{0.30\textwidth} p{0.32\textwidth}}
\toprule
现象 & 可能根因 & 对策 \\
\midrule
计算力矩前馈算出的力矩超出电机量程 & 用了运动学加速度盒、未做力矩可行性检验；或位形导致有效惯量大 & 用 RNEA 在规划期做力矩约束检验（\cref{ins:ad-kin-not-enough}、\cref{ch:arm-traj}）\\
操作空间控制（OSC）末端漂、收不住（$\|\Delta\mathbf{x}\|$ 数十 mm）& 质量阵病态（$\mathrm{cond}(\mathbf{M})\sim10^3$--$10^4$），动力学一致逆数值不稳；正则化破坏投影性 & 换\emph{速度级}冗余解析（\cref{ins:ad-osc-illcond,ins:ad-selection-law}）；别加正则化 \\
$\|\boldsymbol{\tau}_{\text{null}}\|$ 爆炸（上万）& 奇异处 $(\mathbf{J}\mathbf{M}^{-1}\mathbf{J}^\top)^{-1}$ 爆 + 轻腕 $\mathbf{M}^{-1}$ 巨大 & 检查 $\mathrm{cond}(\mathbf{M})$（用真几何惯量）；避奇异位形 \\
低惯量腕关节计算力矩跟踪全幅发散 & 自然频率 $\omega_n=\sqrt{k_p/I}$ 高于控制率，欠采样 & 解耦腕（命令加速度置零 + 物理阻尼兜底），\emph{不}对腕用均匀增益 PD（\cref{ins:ad-wn}）\\
主关节（肩肘）力矩异常（该 $3.5$ 却跳到 $15$）& 漂掉的腕误差经质量阵耦合项 $M_{ij}$ 注入主关节 & 解耦腕的命令加速度（\cref{ins:ad-selection-law} 第 2 条）\\
仿真里 $\mathbf{M}$ “能跑”但换真几何就爆 & 占位惯量是“条件数拐杖”，真几何惯量小 $200\times$ & 用真几何惯量做条件数体检（\cref{ins:ad-placeholder}）\\
科氏项相关的 Lyapunov 证明不消项 & 用了非 Christoffel 的 $\mathbf{C}$，无斜对称 & 改用 Christoffel 选法（\cref{eq:ad-coriolis-matrix}），或先验证 $\dot{\mathbf{M}}-2\mathbf{C}$ 斜对称 \\
\bottomrule
\end{tabular}
\end{center}
